

\section{Related Work}

\noindent\textbf{Video Generation.} The field aims to create temporally coherent and visually realistic dynamic visual content. From early VAE~\cite{vqvae,DBLP:conf/icml/DentonF18,he2018probabilistic,babaeizadeh2021fitvid,4dfy} and GAN~\cite{DBLP:conf/nips/VondrickPT16,tulyakov2018mocogan,clark2019adversarial} frameworks to modern large-scale diffusion models~\cite{VDM,make-a-video,khachatryan2023text2video,wan2025wan,kling} trained on extensive video data, the realism and resolution of generated videos have been revolutionized. In addition to text control generation, researchers have introduced various guidance strategies for video generation, such as structure-based~\cite{ma2024follow,xing2024make,chen2025echomimic}, image-based~\cite{chen2023videocrafter1,deng2024autoregressive}, and temporal controls~\cite{shi2024motion,wu2024motionbooth,yang2024direct}. However, most of them are fundamentally limited as they operate in 2D pixel space, lacking explicit 3D scene modeling. 

\vspace{4pt}
\noindent\textbf{Novel-View Synthesis (NVS).}
3D reconstruction-based methods reconstruct a 3D or 4D representation~\cite{4dgs,zhou2023nerflix,zhu2024fsgs,charatan2024pixelsplat}, ensuring strong geometric consistency but often requiring costly optimization and suffering from artifacts. Generation-based methods use pretrained video diffusion models conditioned on camera trajectories to synthesize novel views~\cite{Reconx,viewcrafter,cat3d,muller2024multidiff,bai2025recammaster}. While leveraging strong generative priors, they often exhibit inconsistencies and object drift under large viewpoint changes or complex scenes. 


\vspace{4pt}
\noindent\textbf{4D Generation.}
4D generation aims to produce temporally evolving 3D representations from texts or sparse images. A growing body of work addresses 3D point tracking~\cite{karaev2024cotracker,karaev2025cotracker3,xiao2024spatialtracker, xiao2025spatialtrackerv2, delta, feng2025st4rtrack, lin2025dgs, lin2025movies} and video-to-4D reconstruction~\cite{xie2024sv4d, wang2025shape, sv4d2, wu2024sc4d, consistent4d, geo4d, yao2025uni4d,wu20254d,gs-dit,cat4d}, estimating dense spatiotemporal correspondences from multi-frame video inputs—a setting where rich temporal signals substantially constrain the problem. Image-to-4D is more ill-posed than video-to-4D, requiring full 3D and temporal reconstruction from single images without multi-frame priors. Existing Image-to-4D methods largely follow two decoupled paradigms: \textit{generate-then-reconstruct}, which first synthesizes videos and then reconstructs 4D (e.g., L4GM~\cite{l4gm}, 4Real~\cite{4real}, DimensionX~\cite{dimensionX}, Free4D~\cite{free4d}), and \textit{reconstruct-then-generate}, which builds static 3D assets before animating them (e.g., Animate124~\cite{animate124}, Animate3D~\cite{animate3d}, 4D-fy~\cite{4dfy}, Gen3C\cite{gen3c}). While effective, the former suffers from video–geometry mismatch, and the latter is often object-centric and limited in scene complexity. Recently, 4DNeX~\cite{4dnex} and One4D~\cite{mi2025one4d} concatenates RGB-XYZ, in contrast to our method, which advances the reconstruct-then-generate paradigm by modeling geometry-aware dense trajectory generation and enables consistent 4D synthesis.

